\section{Picard ranks and fields defining a line}
\label{sec:picard-lines}

\subsection{The geometric Picard lattice}

Over \(\Qbar\), a smooth cubic surface is the blow-up of \(\PP^2\) at six
points.  Write
\[
 \Pic(Y_{\Qbar})
 =\mathbf ZH\oplus\mathbf ZE_1\oplus\cdots\oplus\mathbf ZE_6,
\]
with intersection form
\begin{equation}
 \operatorname{diag}(1,-1,-1,-1,-1,-1,-1).
\label{eq:intersection-form}
\end{equation}
Thus \(\rho(Y_{\Qbar})=7\).  The \(27\) line classes are
\begin{equation}
 E_i,\qquad H-E_i-E_j\ (i<j),\qquad
 2H-\sum_{j\ne i}E_j .
\label{eq:line-classes}
\end{equation}

The six roots
\begin{equation}
 E_1-E_2,\ E_2-E_3,\ E_3-E_4,\ E_4-E_5,\ E_5-E_6,\
 H-E_1-E_2-E_3
\label{eq:simple-roots}
\end{equation}
all have self-intersection \(-2\).  For such a root \(\alpha\), the
reflection is
\begin{equation}
 s_\alpha(x)=x+(x\mathbin{\cdot}\alpha)\alpha.
\label{eq:reflection}
\end{equation}
The exact replay verifies that these six maps preserve
\cref{eq:intersection-form}, the \(27\) classes in
\cref{eq:line-classes}, and their full intersection matrix.  They generate
a group of order \(51840\), act transitively on the line classes, and have
determinant \(-1\) on the \(E_6\) root space.

The anticanonical class
\[
 -K_Y=3H-E_1-\cdots-E_6
\]
is fixed by every reflection.  The six roots span its orthogonal complement.
Therefore a rational Picard vector fixed by all reflections is a multiple of
\(-K_Y\), and
\begin{equation}
 \rank\left(
 \Pic(Y_{\Qbar})\otimes\Q
 \right)^{\WE}=1.
\label{eq:weyl-fixed-rank}
\end{equation}
The independent matrix calculation gives the same fixed rank by stacking
the six matrices \(s_\alpha-I\).

\subsection{Hochschild--Serre at the rank level}

By \cref{thm:weyl}, the Galois image on the geometric Picard lattice is the
full group \(\WE\).  Hence \cref{eq:weyl-fixed-rank} gives
\[
 \rank\Pic(Y_{\Qbar})^{G_{\Q}}=1.
\]
The low-degree Hochschild--Serre sequence, together with Hilbert~90, contains
\begin{equation}
 0\longrightarrow\Pic(Y)
 \longrightarrow\Pic(Y_{\Qbar})^{G_{\Q}}
 \longrightarrow\Br(\Q),
\label{eq:hs-picard}
\end{equation}
as recorded, for example, in Viray, \S2.1, equation~(24) and the following
exact sequence \cite{Viray2023}.  Since \(\Br(\Q)\) is torsion, the cokernel
of the Picard injection is torsion.  Tensoring
\cref{eq:hs-picard} with \(\Q\) therefore gives equality of ranks.
Moreover, \(\operatorname{Pic}^0(Y_{\Qbar})=0\) because a smooth cubic
surface is geometrically rational.  Thus its Picard group equals its
N\'eron--Severi group, and the common rank is the arithmetic Picard number:
\begin{equation}
 \rho(Y/\Q)
 =\rank\Pic(Y)
 =\rank\Pic(Y_{\Qbar})^{G_{\Q}}
 =1.
\label{eq:arithmetic-picard}
\end{equation}
We do not assert an integral equality
\(\Pic(Y)=\Pic(Y_{\Qbar})^{G_{\Q}}\).

\subsection{Line-definition fields}

\begin{corollary}
\label{cor:line-fields}
Let \(L/\Q\) be finite.  If \(Y\) has a line defined over \(L\), then
\[
 27\mid[L:\Q].
\]
In particular \(Y\) has no \(\Q\)-rational line.
\end{corollary}

\begin{proof}
An \(L\)-defined line is an \(L\)-point of \(\FanoY\).  By
\cref{prop:g-irreducible}, this is a \(\Q\)-algebra map \(E\to L\).
Because \(E\) is a field, the map is injective.  Its image is a conjugate
\(E'\subset L\) with \([E':\Q]=27\), and the tower law gives
\[
 [L:\Q]=[L:E'][E':\Q]=27[L:E'].
\]
\end{proof}

\begin{remark}[No rationality inference]
\label{rem:no-rationality}
\Cref{cor:line-fields} concerns rational lines, not rational points.
It proves neither \(Y(\Q)=\varnothing\) nor any statement about rationality,
stable rationality, the Hasse principle, or a Brauer--Manin obstruction.
\end{remark}

\begin{corollary}[Projective invariance]
\label{cor:projective-invariance}
A rational projective coordinate change and multiplication of \(F\) by an
element of \(\Q^\times\) preserve the isomorphism class of \(\FanoY\), the
fields \(E\) and \(K\) up to \(\Q\)-isomorphism, the Galois permutation
representation on the \(27\) lines, and both Picard ranks.
\end{corollary}

\begin{proof}
A coordinate change transports the line section
\cref{eq:line-section} along the induced Grassmannian automorphism, while a
common nonzero scalar leaves its zero scheme unchanged.  The remaining
objects and ranks are functorial consequences of that finite scheme and its
incidence configuration.
\end{proof}
